% MapSafe Mobile results section.
% Required packages in the parent manuscript:
%   \usepackage{graphicx}
%   \usepackage{subcaption}
%   \usepackage{booktabs}
%   \usepackage{tabularx}
%
% The default path assumes that the manuscript is compiled from this directory.
% Define it before % MapSafe Mobile results section.
% Required packages in the parent manuscript:
%   \usepackage{graphicx}
%   \usepackage{subcaption}
%   \usepackage{booktabs}
%   \usepackage{tabularx}
%
% The default path assumes that the manuscript is compiled from this directory.
% Define it before % MapSafe Mobile results section.
% Required packages in the parent manuscript:
%   \usepackage{graphicx}
%   \usepackage{subcaption}
%   \usepackage{booktabs}
%   \usepackage{tabularx}
%
% The default path assumes that the manuscript is compiled from this directory.
% Define it before % MapSafe Mobile results section.
% Required packages in the parent manuscript:
%   \usepackage{graphicx}
%   \usepackage{subcaption}
%   \usepackage{booktabs}
%   \usepackage{tabularx}
%
% The default path assumes that the manuscript is compiled from this directory.
% Define it before \input{mapsafe_results.tex} when overriding the default, e.g.:
%   \newcommand{\mapsafefigdir}{paper/mapsafe-results/screenshots}
\providecommand{\mapsafefigdir}{screenshots}

\section{Results}
\label{sec:mapsafe-results}

The implementation produced an integrated MapSafe Mobile workflow within the
existing NextGIS-based Android application.  The resulting interface supports
two complementary modes.  The guided \emph{Protect \& Share} path presents
spatial representation before encryption, whereas the individual-feature path
allows anonymisation, encryption, and the future notarisation stage to be
entered separately.  Access operations are separated into verification,
decryption, and layer viewing.  Figure~\ref{fig:mapsafe-workflow-entry} shows
these entry points and their relationship.  All nested dialogs and activities
provide a labelled Back control and Android system Back follows the same parent
hierarchy; thus, for example, returning from \emph{Security \& Sharing} leads to
the MapSafe chooser rather than exiting the host map activity.

To make the workflow reproducible when no suitable field layer is available,
the chooser, guided workflow, and anonymisation screen provide a prominent
\emph{Use sample dataset} action (Fig.~\ref{fig:mapsafe-workflow-entry}a--b).
This action creates a known local point-vector layer containing 30 synthetic
observations in six clusters around Suva, selects it, and zooms the map to its
extent.  The sample is created directly as the point geometry type expected by
the anonymisation processors.  Consequently, it also provides a reliable path
past unsupported-file or ``not vector'' errors during demonstrations and
evaluation.

\begin{figure*}[t]
  \centering
  \begin{subfigure}[t]{0.235\textwidth}
    \includegraphics[width=\linewidth]{\mapsafefigdir/paper-01-mapsafe-workflow-chooser.png}
    \caption{MapSafe workflow chooser and sample-data path.}
    \label{fig:mapsafe-workflow-chooser}
  \end{subfigure}\hfill
  \begin{subfigure}[t]{0.235\textwidth}
    \includegraphics[width=\linewidth]{\mapsafefigdir/paper-02-guided-protect-share.png}
    \caption{Guided representation-before-encryption path.}
    \label{fig:mapsafe-guided-workflow}
  \end{subfigure}\hfill
  \begin{subfigure}[t]{0.235\textwidth}
    \includegraphics[width=\linewidth]{\mapsafefigdir/paper-03-safeguard-features.png}
    \caption{Individually selectable safeguard features.}
    \label{fig:mapsafe-safeguard-menu}
  \end{subfigure}\hfill
  \begin{subfigure}[t]{0.235\textwidth}
    \includegraphics[width=\linewidth]{\mapsafefigdir/paper-07-access-features.png}
    \caption{Verification, decryption, and viewing entry points.}
    \label{fig:mapsafe-access-menu}
  \end{subfigure}
  \caption{MapSafe Mobile workflow organisation.  The user may follow the
  ordered Protect \& Share sequence or invoke safeguards individually.  The
  full MapSafe wordmark identifies the main chooser, while compact headers use
  the icon-only mark.}
  \label{fig:mapsafe-workflow-entry}
\end{figure*}

\subsection{Anonymisation and spatial representation}
\label{subsec:mapsafe-anonymisation-results}

The anonymisation branch offers halo (donut) masking and hexagonal binning
(Fig.~\ref{fig:mapsafe-anonymisation-controls}).  Halo masking preserves every
feature and its non-spatial attributes but replaces each WGS84 point
$p_i=(\lambda_i,\phi_i)$ with a masked point $q_i$.  For each feature, a
cryptographically secure random generator selects a bearing uniformly from
$[0,2\pi)$ and a displacement uniformly from the user-selected interval
$[r_{\min},r_{\max})$.  The destination is calculated on a sphere rather than
by adding degree offsets.  The source layer remains unchanged and the candidate
masked layer is displayed in blue, allowing it to be compared with the precise
red reference points.

The implementation calculates Spruill's disclosure-risk measure immediately
after every masking attempt.  For $n$ paired original and masked points, define

\begin{equation}
 I_i =
 \begin{cases}
 1, & d(q_i,p_i)=\min_{1\leq j\leq n} d(q_i,p_j),\\
 0, & \text{otherwise},
 \end{cases}
 \qquad
 R_{\mathrm{Spruill}}=\frac{100}{n}\sum_{i=1}^{n}I_i .
 \label{eq:spruill-risk}
\end{equation}

Here, a point contributes to disclosure risk when the closest original location
to its masked position is still its own parent.  Nearest-neighbour searches are
performed with a three-dimensional unit-sphere $k$-d tree; unit-sphere chord
distance is monotonic with great-circle distance and avoids latitude-dependent
planar distortion.  Equal-distance ties that include the parent are counted as
parent-nearest.  Because a higher value is easier to interpret as stronger
privacy, the interface reports the inverted measure

\begin{equation}
 P_{\mathrm{inv}} = 100-R_{\mathrm{Spruill}}.
 \label{eq:spruill-inverted}
\end{equation}

In the illustrative 30-point run captured in
Figure~\ref{fig:mapsafe-anonymisation-result}a, 4 points remained
parent-nearest.  This gave $R_{\mathrm{Spruill}}=13.33\%$ and
$P_{\mathrm{inv}}=86.67/100$, with a mean displacement of 1041.21~m for the
selected 100--2000~m halo.  These values describe this stochastic run rather
than a fixed benchmark result.  The result screen exposes both the aggregate
score and its component counts.  The user may accept the candidate or select
\emph{Remask}; importantly, remasking returns to the same controls but starts
again from the precise source layer, never from the already displaced
candidate (Fig.~\ref{fig:mapsafe-anonymisation-result}b).  This prevents
unintended compounding of random offsets.

Hexagonal binning removes the individual point representation.  Each point is
assigned to a cell at the selected resolution, points with the same cell ID are
grouped, and one polygon containing a count attribute is written for each
occupied cell.  Native ARM and ARM64 builds use H3.  The x86/x86\_64 emulator
used for these captures lacks the H3 native library and therefore selected the
native-free Web Mercator hexagonal grid, whose resolution-dependent cell sizes
approximate the corresponding H3 choices.  The protected map in
Figure~\ref{fig:mapsafe-anonymisation-result}c demonstrates the resulting
masked points and counted polygon representation over the Suva sample.

\begin{figure*}[t]
  \centering
  \begin{subfigure}[t]{0.31\textwidth}
    \includegraphics[width=\linewidth]{\mapsafefigdir/paper-04-anonymise-options.png}
    \caption{Choice of masking or aggregation.}
    \label{fig:mapsafe-anonymise-choice}
  \end{subfigure}\hfill
  \begin{subfigure}[t]{0.31\textwidth}
    \includegraphics[width=\linewidth]{\mapsafefigdir/paper-16-halo-configuration.png}
    \caption{Halo bounds and privacy control.}
    \label{fig:mapsafe-halo-controls}
  \end{subfigure}\hfill
  \begin{subfigure}[t]{0.31\textwidth}
    \includegraphics[width=\linewidth]{\mapsafefigdir/paper-19-hexbin-configuration.png}
    \caption{Hexagon resolution and count aggregation.}
    \label{fig:mapsafe-hexbin-controls}
  \end{subfigure}
  \caption{Anonymisation controls.  Halo masking retains one displaced feature
  per input feature, whereas hexagonal binning replaces locations with occupied
  cells and their counts.}
  \label{fig:mapsafe-anonymisation-controls}
\end{figure*}

\begin{figure*}[t]
  \centering
  \begin{subfigure}[t]{0.31\textwidth}
    \includegraphics[width=\linewidth]{\mapsafefigdir/paper-17-halo-spruill-result.png}
    \caption{Inverted Spruill score and component measures.}
    \label{fig:mapsafe-spruill-result}
  \end{subfigure}\hfill
  \begin{subfigure}[t]{0.31\textwidth}
    \includegraphics[width=\linewidth]{\mapsafefigdir/paper-18-halo-remask.png}
    \caption{Remasking explicitly restarts from the precise source.}
    \label{fig:mapsafe-remask}
  \end{subfigure}\hfill
  \begin{subfigure}[t]{0.31\textwidth}
    \includegraphics[width=\linewidth]{\mapsafefigdir/paper-20-protected-map.png}
    \caption{Map output with precise references, masked points, and hexagons.}
    \label{fig:mapsafe-protected-map}
  \end{subfigure}
  \caption{Anonymisation outcomes for the bundled Suva point dataset.  The
  numeric privacy result is shown before the candidate is accepted, preserving
  an opportunity to remask or change parameters.}
  \label{fig:mapsafe-anonymisation-result}
\end{figure*}

\subsection{Encryption identity and public-key exchange}
\label{subsec:mapsafe-key-results}

The public-key workflow is surfaced through \emph{Security \& Sharing}
(Fig.~\ref{fig:mapsafe-key-exchange}).  Its detailed sequence is as follows.

\begin{enumerate}
  \item The user creates a local OpenPGP identity from a name, organisation,
  optional email address, and a passphrase of at least 12 characters
  (Fig.~\ref{fig:mapsafe-key-exchange}a).  The device generates an RSA-3072
  primary key for certification and signing plus a separate RSA-3072 encryption
  subkey.  The resulting full fingerprint is displayed to the user
  (Fig.~\ref{fig:mapsafe-key-exchange}b).

  \item The secret key material is encrypted with the user's passphrase and is
  additionally wrapped at rest by a non-exportable AES key in Android Keystore.
  It is not uploaded during ordinary key exchange.  Only the public key is
  exportable or publishable; an explicitly protected private-key backup is
  available as a recovery operation.

  \item The user selects an authenticated NextGIS account and a trusted NextGIS
  group (Fig.~\ref{fig:mapsafe-key-exchange}c).  MapSafe publishes the local
  public key to a group-specific key directory and can download keys belonging
  to current group members (Fig.~\ref{fig:mapsafe-key-exchange}d).  Directory
  ownership and current authentication-group membership are checked during
  synchronisation.

  \item Downloaded keys are initially discovered, not silently trusted.  The
  interface instructs the user to compare the complete fingerprint through an
  independent trusted channel and then accept it explicitly.  Acceptance pins
  that fingerprint to the member.  A later key change enters a pending/quarantined
  state, while changed, missing, duplicate, revoked, or removed-member keys are
  blocked from new recipient selection.  Previously accepted public keys remain
  cached so that encryption for established recipients can proceed offline.
\end{enumerate}

This design uses the authenticated group to establish membership and distribute
public material, but does not equate group membership with cryptographic key
authenticity.  The explicit fingerprint step supplies that missing binding and
prevents an updated directory entry from silently replacing a recipient's key.

\begin{figure*}[t]
  \centering
  \begin{subfigure}[t]{0.235\textwidth}
    \includegraphics[width=\linewidth]{\mapsafefigdir/paper-10-identity-creation.png}
    \caption{Local identity and passphrase entry.}
  \end{subfigure}\hfill
  \begin{subfigure}[t]{0.235\textwidth}
    \includegraphics[width=\linewidth]{\mapsafefigdir/paper-11-identity-success.png}
    \caption{Generated identity, key ID, and fingerprint.}
  \end{subfigure}\hfill
  \begin{subfigure}[t]{0.235\textwidth}
    \includegraphics[width=\linewidth]{\mapsafefigdir/paper-14-security-account-group.png}
    \caption{NextGIS account and group selection.}
  \end{subfigure}\hfill
  \begin{subfigure}[t]{0.235\textwidth}
    \includegraphics[width=\linewidth]{\mapsafefigdir/paper-15-security-identity-keys.png}
    \caption{Publishing, synchronising, and reviewing group keys.}
  \end{subfigure}
  \caption{Encryption identity and public-key exchange.  Private key material
  remains locally protected; public keys are distributed within a selected
  NextGIS group and become eligible recipients only after fingerprint review.}
  \label{fig:mapsafe-key-exchange}
\end{figure*}

\subsection{Multi-recipient encryption and signing}
\label{subsec:mapsafe-encryption-results}

Encryption may receive the currently selected protected map layer directly or
an existing file from device storage (Fig.~\ref{fig:mapsafe-encryption-flow}a).
For a map layer, MapSafe first exports the active representation as WGS84
GeoJSON, so the precise input need not leave the device.  The encryption screen
then shows the local identity, accepted recipient count, and the selected
content-protection profile (Fig.~\ref{fig:mapsafe-encryption-flow}b).

The package uses the RFC~9580 version-2 Symmetrically Encrypted Integrity
Protected Data profile with AES-256-GCM.  A fresh random content-encryption key
encrypts the compressed file payload once.  For each checked recipient, OpenPGP
adds a public-key encrypted session-key packet: the recipient's public
encryption subkey protects the content key, not a second copy of the entire
dataset.  Therefore adding recipients increases only the key-envelope portion
of the package.  The device owner's own public key can be selected as a
self-recipient, ensuring that the sender can subsequently decrypt the package.
The selection dialog makes this relationship explicit
(Fig.~\ref{fig:mapsafe-encryption-flow}c).

When \emph{Sign with my local identity} is selected, the passphrase unlocks the
local signing key only for the operation and a SHA-256 OpenPGP binary-document
signature is embedded with the payload.  Encryption provides confidentiality
and AEAD integrity, whereas this signature separately allows the receiver to
test which key created the package.  The successful state reports the
AES-256-GCM profile, authorised-recipient count, and whether a signed package
was saved (Fig.~\ref{fig:mapsafe-encryption-flow}d).

\begin{figure*}[t]
  \centering
  \begin{subfigure}[t]{0.235\textwidth}
    \includegraphics[width=\linewidth]{\mapsafefigdir/paper-05-encrypt-source-chooser.png}
    \caption{Protected-layer or file source selection.}
  \end{subfigure}\hfill
  \begin{subfigure}[t]{0.235\textwidth}
    \includegraphics[width=\linewidth]{\mapsafefigdir/paper-21-encrypt-protect.png}
    \caption{Encryption identity and protection summary.}
  \end{subfigure}\hfill
  \begin{subfigure}[t]{0.235\textwidth}
    \includegraphics[width=\linewidth]{\mapsafefigdir/paper-12-recipient-selection.png}
    \caption{Verified recipients and optional signing.}
  \end{subfigure}\hfill
  \begin{subfigure}[t]{0.235\textwidth}
    \includegraphics[width=\linewidth]{\mapsafefigdir/paper-22-encryption-success.png}
    \caption{Signed package ready to share.}
  \end{subfigure}
  \caption{Multi-recipient encryption flow.  The payload is encrypted once with
  AES-256-GCM; a separately public-key-encrypted session key grants access to
  each selected recipient.}
  \label{fig:mapsafe-encryption-flow}
\end{figure*}

\subsection{Decryption, verification, and viewing}
\label{subsec:mapsafe-decryption-results}

The access workflow starts with an OpenPGP package selected from storage
(Fig.~\ref{fig:mapsafe-decryption-flow}a--b).  MapSafe inspects the package's
recipient identifiers and chooses a matching locally held encryption secret
key.  The user passphrase unlocks that key; the resulting private key decrypts
the recipient's session-key packet, after which AES-256-GCM decrypts the single
payload.  A package addressed only to other recipients is rejected before
plaintext is released.

Two independent checks follow.  First, the SEIPD/AEAD integrity result must be
valid; an unprotected or modified package is rejected.  Second, if an embedded
signature is present, MapSafe locates the corresponding verification public key
and reports the signature as valid, invalid, or from an unknown signer.  The
decrypted file is held as a candidate until these checks complete.  In the
tested path the interface reported both a passed integrity check and a
cryptographically valid signature before offering \emph{Import and open}
(Fig.~\ref{fig:mapsafe-decryption-flow}c).  The final screen retains this status
as visible evidence (Fig.~\ref{fig:mapsafe-decryption-flow}d).

For a valid GeoJSON result, choosing \emph{Import and open} creates a local
NextGIS vector layer, selects it, and zooms to its extent.  Invalidly signed data
is not offered for map import.  This final promotion step separates successful
cryptographic processing from viewing and prevents an unauthenticated temporary
output from being mistaken for an accepted layer.

\begin{figure*}[t]
  \centering
  \begin{subfigure}[t]{0.235\textwidth}
    \includegraphics[width=\linewidth]{\mapsafefigdir/paper-08-decrypt-package-chooser.png}
    \caption{OpenPGP package entry point.}
  \end{subfigure}\hfill
  \begin{subfigure}[t]{0.235\textwidth}
    \includegraphics[width=\linewidth]{\mapsafefigdir/paper-23-decrypt-access.png}
    \caption{Local private-key and integrity summary.}
  \end{subfigure}\hfill
  \begin{subfigure}[t]{0.235\textwidth}
    \includegraphics[width=\linewidth]{\mapsafefigdir/paper-13-decryption-integrity-prompt.png}
    \caption{Verified candidate offered for import.}
  \end{subfigure}\hfill
  \begin{subfigure}[t]{0.235\textwidth}
    \includegraphics[width=\linewidth]{\mapsafefigdir/paper-24-decryption-success.png}
    \caption{Persistent success and signature status.}
  \end{subfigure}
  \caption{Decryption and access flow.  MapSafe first proves recipient access,
  then verifies AEAD integrity and the optional OpenPGP signature, and only then
  offers a valid spatial payload for import and viewing.}
  \label{fig:mapsafe-decryption-flow}
\end{figure*}

\subsection{Notarisation and blockchain-verification boundary}
\label{subsec:mapsafe-blockchain-results}

The proposed notarisation design records a SHA-256 digest of the encrypted
package plus minimal metadata, while the encrypted dataset remains off-chain.
Verification would independently hash the candidate package and compare that
digest with the value associated with a blockchain receipt; equality would
confirm byte-level correspondence with the notarised artifact, while inequality
would identify modification or selection of the wrong package.  This check is
complementary to OpenPGP: blockchain comparison provides timestamped external
evidence for a package digest, whereas decryption, AEAD integrity, and signature
verification establish authorised access and package authenticity.

However, the present Android build does \emph{not} contain a production
blockchain client and did not submit a transaction.  The interface deliberately
states that the production network is not configured and disables the minting
action (Fig.~\ref{fig:mapsafe-blockchain-boundary}a); the verification menu is
likewise marked as a later-stage interface
(Fig.~\ref{fig:mapsafe-blockchain-boundary}b).  This is an important result
boundary: these screens demonstrate workflow placement and privacy constraints,
not completed on-chain notarisation.  Automated scenarios currently exercise
the SHA-256 receipt comparison with an in-memory test ledger, including rejection
of a modified artifact, but those tests are not evidence of deployment to a
real network.

\begin{figure*}[t]
  \centering
  \begin{subfigure}[t]{0.31\textwidth}
    \includegraphics[width=\linewidth]{\mapsafefigdir/paper-06-notarisation-placeholder.png}
    \caption{Notarisation screen with the production action disabled.}
  \end{subfigure}\hspace{0.05\textwidth}
  \begin{subfigure}[t]{0.31\textwidth}
    \includegraphics[width=\linewidth]{\mapsafefigdir/paper-09-verification-placeholder.png}
    \caption{Blockchain-verification interface identified as future work.}
  \end{subfigure}
  \caption{Current blockchain boundary.  No real blockchain transaction is
  simulated or presented as completed by this build.}
  \label{fig:mapsafe-blockchain-boundary}
\end{figure*}

\subsection{Automated validation}
\label{subsec:mapsafe-validation-results}

The host-side MapSafe test suite executed 27 tests with no failures or skipped
tests.  These comprised eight OpenPGP engine tests, six public-key trust-state
tests, five Spruill-measure tests, two donut-masking tests, two portable-hex-grid
tests, and four end-to-end workflow-scenario tests.  Their coverage included
multi-recipient encryption, signed and unsigned packages, integrity/tamper
rejection, recipient mismatch, fingerprint changes, workflow combinations, and
the in-memory digest ledger.

Five targeted instrumented tests were then executed on a Pixel~9a Android
Virtual Device, again with no failures or skipped tests.  They traversed and
captured the workflow hierarchy, generated an actual RSA-3072 identity, created
the bundled point-vector layer, performed masking and hexagonal binning, exposed
the Security \& Sharing key-exchange journey, and completed a real UI-driven
signed AES-256-GCM encrypt--decrypt round trip.  The 24 screen states reproduced
in Figures~\ref{fig:mapsafe-workflow-entry}--\ref{fig:mapsafe-blockchain-boundary}
were extracted from these instrumented workflows.  Thus, the screenshots are
not design mock-ups; they are outputs of the implemented Android application.
 when overriding the default, e.g.:
%   \newcommand{\mapsafefigdir}{paper/mapsafe-results/screenshots}
\providecommand{\mapsafefigdir}{screenshots}

\section{Results}
\label{sec:mapsafe-results}

The implementation produced an integrated MapSafe Mobile workflow within the
existing NextGIS-based Android application.  The resulting interface supports
two complementary modes.  The guided \emph{Protect \& Share} path presents
spatial representation before encryption, whereas the individual-feature path
allows anonymisation, encryption, and the future notarisation stage to be
entered separately.  Access operations are separated into verification,
decryption, and layer viewing.  Figure~\ref{fig:mapsafe-workflow-entry} shows
these entry points and their relationship.  All nested dialogs and activities
provide a labelled Back control and Android system Back follows the same parent
hierarchy; thus, for example, returning from \emph{Security \& Sharing} leads to
the MapSafe chooser rather than exiting the host map activity.

To make the workflow reproducible when no suitable field layer is available,
the chooser, guided workflow, and anonymisation screen provide a prominent
\emph{Use sample dataset} action (Fig.~\ref{fig:mapsafe-workflow-entry}a--b).
This action creates a known local point-vector layer containing 30 synthetic
observations in six clusters around Suva, selects it, and zooms the map to its
extent.  The sample is created directly as the point geometry type expected by
the anonymisation processors.  Consequently, it also provides a reliable path
past unsupported-file or ``not vector'' errors during demonstrations and
evaluation.

\begin{figure*}[t]
  \centering
  \begin{subfigure}[t]{0.235\textwidth}
    \includegraphics[width=\linewidth]{\mapsafefigdir/paper-01-mapsafe-workflow-chooser.png}
    \caption{MapSafe workflow chooser and sample-data path.}
    \label{fig:mapsafe-workflow-chooser}
  \end{subfigure}\hfill
  \begin{subfigure}[t]{0.235\textwidth}
    \includegraphics[width=\linewidth]{\mapsafefigdir/paper-02-guided-protect-share.png}
    \caption{Guided representation-before-encryption path.}
    \label{fig:mapsafe-guided-workflow}
  \end{subfigure}\hfill
  \begin{subfigure}[t]{0.235\textwidth}
    \includegraphics[width=\linewidth]{\mapsafefigdir/paper-03-safeguard-features.png}
    \caption{Individually selectable safeguard features.}
    \label{fig:mapsafe-safeguard-menu}
  \end{subfigure}\hfill
  \begin{subfigure}[t]{0.235\textwidth}
    \includegraphics[width=\linewidth]{\mapsafefigdir/paper-07-access-features.png}
    \caption{Verification, decryption, and viewing entry points.}
    \label{fig:mapsafe-access-menu}
  \end{subfigure}
  \caption{MapSafe Mobile workflow organisation.  The user may follow the
  ordered Protect \& Share sequence or invoke safeguards individually.  The
  full MapSafe wordmark identifies the main chooser, while compact headers use
  the icon-only mark.}
  \label{fig:mapsafe-workflow-entry}
\end{figure*}

\subsection{Anonymisation and spatial representation}
\label{subsec:mapsafe-anonymisation-results}

The anonymisation branch offers halo (donut) masking and hexagonal binning
(Fig.~\ref{fig:mapsafe-anonymisation-controls}).  Halo masking preserves every
feature and its non-spatial attributes but replaces each WGS84 point
$p_i=(\lambda_i,\phi_i)$ with a masked point $q_i$.  For each feature, a
cryptographically secure random generator selects a bearing uniformly from
$[0,2\pi)$ and a displacement uniformly from the user-selected interval
$[r_{\min},r_{\max})$.  The destination is calculated on a sphere rather than
by adding degree offsets.  The source layer remains unchanged and the candidate
masked layer is displayed in blue, allowing it to be compared with the precise
red reference points.

The implementation calculates Spruill's disclosure-risk measure immediately
after every masking attempt.  For $n$ paired original and masked points, define

\begin{equation}
 I_i =
 \begin{cases}
 1, & d(q_i,p_i)=\min_{1\leq j\leq n} d(q_i,p_j),\\
 0, & \text{otherwise},
 \end{cases}
 \qquad
 R_{\mathrm{Spruill}}=\frac{100}{n}\sum_{i=1}^{n}I_i .
 \label{eq:spruill-risk}
\end{equation}

Here, a point contributes to disclosure risk when the closest original location
to its masked position is still its own parent.  Nearest-neighbour searches are
performed with a three-dimensional unit-sphere $k$-d tree; unit-sphere chord
distance is monotonic with great-circle distance and avoids latitude-dependent
planar distortion.  Equal-distance ties that include the parent are counted as
parent-nearest.  Because a higher value is easier to interpret as stronger
privacy, the interface reports the inverted measure

\begin{equation}
 P_{\mathrm{inv}} = 100-R_{\mathrm{Spruill}}.
 \label{eq:spruill-inverted}
\end{equation}

In the illustrative 30-point run captured in
Figure~\ref{fig:mapsafe-anonymisation-result}a, 4 points remained
parent-nearest.  This gave $R_{\mathrm{Spruill}}=13.33\%$ and
$P_{\mathrm{inv}}=86.67/100$, with a mean displacement of 1041.21~m for the
selected 100--2000~m halo.  These values describe this stochastic run rather
than a fixed benchmark result.  The result screen exposes both the aggregate
score and its component counts.  The user may accept the candidate or select
\emph{Remask}; importantly, remasking returns to the same controls but starts
again from the precise source layer, never from the already displaced
candidate (Fig.~\ref{fig:mapsafe-anonymisation-result}b).  This prevents
unintended compounding of random offsets.

Hexagonal binning removes the individual point representation.  Each point is
assigned to a cell at the selected resolution, points with the same cell ID are
grouped, and one polygon containing a count attribute is written for each
occupied cell.  Native ARM and ARM64 builds use H3.  The x86/x86\_64 emulator
used for these captures lacks the H3 native library and therefore selected the
native-free Web Mercator hexagonal grid, whose resolution-dependent cell sizes
approximate the corresponding H3 choices.  The protected map in
Figure~\ref{fig:mapsafe-anonymisation-result}c demonstrates the resulting
masked points and counted polygon representation over the Suva sample.

\begin{figure*}[t]
  \centering
  \begin{subfigure}[t]{0.31\textwidth}
    \includegraphics[width=\linewidth]{\mapsafefigdir/paper-04-anonymise-options.png}
    \caption{Choice of masking or aggregation.}
    \label{fig:mapsafe-anonymise-choice}
  \end{subfigure}\hfill
  \begin{subfigure}[t]{0.31\textwidth}
    \includegraphics[width=\linewidth]{\mapsafefigdir/paper-16-halo-configuration.png}
    \caption{Halo bounds and privacy control.}
    \label{fig:mapsafe-halo-controls}
  \end{subfigure}\hfill
  \begin{subfigure}[t]{0.31\textwidth}
    \includegraphics[width=\linewidth]{\mapsafefigdir/paper-19-hexbin-configuration.png}
    \caption{Hexagon resolution and count aggregation.}
    \label{fig:mapsafe-hexbin-controls}
  \end{subfigure}
  \caption{Anonymisation controls.  Halo masking retains one displaced feature
  per input feature, whereas hexagonal binning replaces locations with occupied
  cells and their counts.}
  \label{fig:mapsafe-anonymisation-controls}
\end{figure*}

\begin{figure*}[t]
  \centering
  \begin{subfigure}[t]{0.31\textwidth}
    \includegraphics[width=\linewidth]{\mapsafefigdir/paper-17-halo-spruill-result.png}
    \caption{Inverted Spruill score and component measures.}
    \label{fig:mapsafe-spruill-result}
  \end{subfigure}\hfill
  \begin{subfigure}[t]{0.31\textwidth}
    \includegraphics[width=\linewidth]{\mapsafefigdir/paper-18-halo-remask.png}
    \caption{Remasking explicitly restarts from the precise source.}
    \label{fig:mapsafe-remask}
  \end{subfigure}\hfill
  \begin{subfigure}[t]{0.31\textwidth}
    \includegraphics[width=\linewidth]{\mapsafefigdir/paper-20-protected-map.png}
    \caption{Map output with precise references, masked points, and hexagons.}
    \label{fig:mapsafe-protected-map}
  \end{subfigure}
  \caption{Anonymisation outcomes for the bundled Suva point dataset.  The
  numeric privacy result is shown before the candidate is accepted, preserving
  an opportunity to remask or change parameters.}
  \label{fig:mapsafe-anonymisation-result}
\end{figure*}

\subsection{Encryption identity and public-key exchange}
\label{subsec:mapsafe-key-results}

The public-key workflow is surfaced through \emph{Security \& Sharing}
(Fig.~\ref{fig:mapsafe-key-exchange}).  Its detailed sequence is as follows.

\begin{enumerate}
  \item The user creates a local OpenPGP identity from a name, organisation,
  optional email address, and a passphrase of at least 12 characters
  (Fig.~\ref{fig:mapsafe-key-exchange}a).  The device generates an RSA-3072
  primary key for certification and signing plus a separate RSA-3072 encryption
  subkey.  The resulting full fingerprint is displayed to the user
  (Fig.~\ref{fig:mapsafe-key-exchange}b).

  \item The secret key material is encrypted with the user's passphrase and is
  additionally wrapped at rest by a non-exportable AES key in Android Keystore.
  It is not uploaded during ordinary key exchange.  Only the public key is
  exportable or publishable; an explicitly protected private-key backup is
  available as a recovery operation.

  \item The user selects an authenticated NextGIS account and a trusted NextGIS
  group (Fig.~\ref{fig:mapsafe-key-exchange}c).  MapSafe publishes the local
  public key to a group-specific key directory and can download keys belonging
  to current group members (Fig.~\ref{fig:mapsafe-key-exchange}d).  Directory
  ownership and current authentication-group membership are checked during
  synchronisation.

  \item Downloaded keys are initially discovered, not silently trusted.  The
  interface instructs the user to compare the complete fingerprint through an
  independent trusted channel and then accept it explicitly.  Acceptance pins
  that fingerprint to the member.  A later key change enters a pending/quarantined
  state, while changed, missing, duplicate, revoked, or removed-member keys are
  blocked from new recipient selection.  Previously accepted public keys remain
  cached so that encryption for established recipients can proceed offline.
\end{enumerate}

This design uses the authenticated group to establish membership and distribute
public material, but does not equate group membership with cryptographic key
authenticity.  The explicit fingerprint step supplies that missing binding and
prevents an updated directory entry from silently replacing a recipient's key.

\begin{figure*}[t]
  \centering
  \begin{subfigure}[t]{0.235\textwidth}
    \includegraphics[width=\linewidth]{\mapsafefigdir/paper-10-identity-creation.png}
    \caption{Local identity and passphrase entry.}
  \end{subfigure}\hfill
  \begin{subfigure}[t]{0.235\textwidth}
    \includegraphics[width=\linewidth]{\mapsafefigdir/paper-11-identity-success.png}
    \caption{Generated identity, key ID, and fingerprint.}
  \end{subfigure}\hfill
  \begin{subfigure}[t]{0.235\textwidth}
    \includegraphics[width=\linewidth]{\mapsafefigdir/paper-14-security-account-group.png}
    \caption{NextGIS account and group selection.}
  \end{subfigure}\hfill
  \begin{subfigure}[t]{0.235\textwidth}
    \includegraphics[width=\linewidth]{\mapsafefigdir/paper-15-security-identity-keys.png}
    \caption{Publishing, synchronising, and reviewing group keys.}
  \end{subfigure}
  \caption{Encryption identity and public-key exchange.  Private key material
  remains locally protected; public keys are distributed within a selected
  NextGIS group and become eligible recipients only after fingerprint review.}
  \label{fig:mapsafe-key-exchange}
\end{figure*}

\subsection{Multi-recipient encryption and signing}
\label{subsec:mapsafe-encryption-results}

Encryption may receive the currently selected protected map layer directly or
an existing file from device storage (Fig.~\ref{fig:mapsafe-encryption-flow}a).
For a map layer, MapSafe first exports the active representation as WGS84
GeoJSON, so the precise input need not leave the device.  The encryption screen
then shows the local identity, accepted recipient count, and the selected
content-protection profile (Fig.~\ref{fig:mapsafe-encryption-flow}b).

The package uses the RFC~9580 version-2 Symmetrically Encrypted Integrity
Protected Data profile with AES-256-GCM.  A fresh random content-encryption key
encrypts the compressed file payload once.  For each checked recipient, OpenPGP
adds a public-key encrypted session-key packet: the recipient's public
encryption subkey protects the content key, not a second copy of the entire
dataset.  Therefore adding recipients increases only the key-envelope portion
of the package.  The device owner's own public key can be selected as a
self-recipient, ensuring that the sender can subsequently decrypt the package.
The selection dialog makes this relationship explicit
(Fig.~\ref{fig:mapsafe-encryption-flow}c).

When \emph{Sign with my local identity} is selected, the passphrase unlocks the
local signing key only for the operation and a SHA-256 OpenPGP binary-document
signature is embedded with the payload.  Encryption provides confidentiality
and AEAD integrity, whereas this signature separately allows the receiver to
test which key created the package.  The successful state reports the
AES-256-GCM profile, authorised-recipient count, and whether a signed package
was saved (Fig.~\ref{fig:mapsafe-encryption-flow}d).

\begin{figure*}[t]
  \centering
  \begin{subfigure}[t]{0.235\textwidth}
    \includegraphics[width=\linewidth]{\mapsafefigdir/paper-05-encrypt-source-chooser.png}
    \caption{Protected-layer or file source selection.}
  \end{subfigure}\hfill
  \begin{subfigure}[t]{0.235\textwidth}
    \includegraphics[width=\linewidth]{\mapsafefigdir/paper-21-encrypt-protect.png}
    \caption{Encryption identity and protection summary.}
  \end{subfigure}\hfill
  \begin{subfigure}[t]{0.235\textwidth}
    \includegraphics[width=\linewidth]{\mapsafefigdir/paper-12-recipient-selection.png}
    \caption{Verified recipients and optional signing.}
  \end{subfigure}\hfill
  \begin{subfigure}[t]{0.235\textwidth}
    \includegraphics[width=\linewidth]{\mapsafefigdir/paper-22-encryption-success.png}
    \caption{Signed package ready to share.}
  \end{subfigure}
  \caption{Multi-recipient encryption flow.  The payload is encrypted once with
  AES-256-GCM; a separately public-key-encrypted session key grants access to
  each selected recipient.}
  \label{fig:mapsafe-encryption-flow}
\end{figure*}

\subsection{Decryption, verification, and viewing}
\label{subsec:mapsafe-decryption-results}

The access workflow starts with an OpenPGP package selected from storage
(Fig.~\ref{fig:mapsafe-decryption-flow}a--b).  MapSafe inspects the package's
recipient identifiers and chooses a matching locally held encryption secret
key.  The user passphrase unlocks that key; the resulting private key decrypts
the recipient's session-key packet, after which AES-256-GCM decrypts the single
payload.  A package addressed only to other recipients is rejected before
plaintext is released.

Two independent checks follow.  First, the SEIPD/AEAD integrity result must be
valid; an unprotected or modified package is rejected.  Second, if an embedded
signature is present, MapSafe locates the corresponding verification public key
and reports the signature as valid, invalid, or from an unknown signer.  The
decrypted file is held as a candidate until these checks complete.  In the
tested path the interface reported both a passed integrity check and a
cryptographically valid signature before offering \emph{Import and open}
(Fig.~\ref{fig:mapsafe-decryption-flow}c).  The final screen retains this status
as visible evidence (Fig.~\ref{fig:mapsafe-decryption-flow}d).

For a valid GeoJSON result, choosing \emph{Import and open} creates a local
NextGIS vector layer, selects it, and zooms to its extent.  Invalidly signed data
is not offered for map import.  This final promotion step separates successful
cryptographic processing from viewing and prevents an unauthenticated temporary
output from being mistaken for an accepted layer.

\begin{figure*}[t]
  \centering
  \begin{subfigure}[t]{0.235\textwidth}
    \includegraphics[width=\linewidth]{\mapsafefigdir/paper-08-decrypt-package-chooser.png}
    \caption{OpenPGP package entry point.}
  \end{subfigure}\hfill
  \begin{subfigure}[t]{0.235\textwidth}
    \includegraphics[width=\linewidth]{\mapsafefigdir/paper-23-decrypt-access.png}
    \caption{Local private-key and integrity summary.}
  \end{subfigure}\hfill
  \begin{subfigure}[t]{0.235\textwidth}
    \includegraphics[width=\linewidth]{\mapsafefigdir/paper-13-decryption-integrity-prompt.png}
    \caption{Verified candidate offered for import.}
  \end{subfigure}\hfill
  \begin{subfigure}[t]{0.235\textwidth}
    \includegraphics[width=\linewidth]{\mapsafefigdir/paper-24-decryption-success.png}
    \caption{Persistent success and signature status.}
  \end{subfigure}
  \caption{Decryption and access flow.  MapSafe first proves recipient access,
  then verifies AEAD integrity and the optional OpenPGP signature, and only then
  offers a valid spatial payload for import and viewing.}
  \label{fig:mapsafe-decryption-flow}
\end{figure*}

\subsection{Notarisation and blockchain-verification boundary}
\label{subsec:mapsafe-blockchain-results}

The proposed notarisation design records a SHA-256 digest of the encrypted
package plus minimal metadata, while the encrypted dataset remains off-chain.
Verification would independently hash the candidate package and compare that
digest with the value associated with a blockchain receipt; equality would
confirm byte-level correspondence with the notarised artifact, while inequality
would identify modification or selection of the wrong package.  This check is
complementary to OpenPGP: blockchain comparison provides timestamped external
evidence for a package digest, whereas decryption, AEAD integrity, and signature
verification establish authorised access and package authenticity.

However, the present Android build does \emph{not} contain a production
blockchain client and did not submit a transaction.  The interface deliberately
states that the production network is not configured and disables the minting
action (Fig.~\ref{fig:mapsafe-blockchain-boundary}a); the verification menu is
likewise marked as a later-stage interface
(Fig.~\ref{fig:mapsafe-blockchain-boundary}b).  This is an important result
boundary: these screens demonstrate workflow placement and privacy constraints,
not completed on-chain notarisation.  Automated scenarios currently exercise
the SHA-256 receipt comparison with an in-memory test ledger, including rejection
of a modified artifact, but those tests are not evidence of deployment to a
real network.

\begin{figure*}[t]
  \centering
  \begin{subfigure}[t]{0.31\textwidth}
    \includegraphics[width=\linewidth]{\mapsafefigdir/paper-06-notarisation-placeholder.png}
    \caption{Notarisation screen with the production action disabled.}
  \end{subfigure}\hspace{0.05\textwidth}
  \begin{subfigure}[t]{0.31\textwidth}
    \includegraphics[width=\linewidth]{\mapsafefigdir/paper-09-verification-placeholder.png}
    \caption{Blockchain-verification interface identified as future work.}
  \end{subfigure}
  \caption{Current blockchain boundary.  No real blockchain transaction is
  simulated or presented as completed by this build.}
  \label{fig:mapsafe-blockchain-boundary}
\end{figure*}

\subsection{Automated validation}
\label{subsec:mapsafe-validation-results}

The host-side MapSafe test suite executed 27 tests with no failures or skipped
tests.  These comprised eight OpenPGP engine tests, six public-key trust-state
tests, five Spruill-measure tests, two donut-masking tests, two portable-hex-grid
tests, and four end-to-end workflow-scenario tests.  Their coverage included
multi-recipient encryption, signed and unsigned packages, integrity/tamper
rejection, recipient mismatch, fingerprint changes, workflow combinations, and
the in-memory digest ledger.

Five targeted instrumented tests were then executed on a Pixel~9a Android
Virtual Device, again with no failures or skipped tests.  They traversed and
captured the workflow hierarchy, generated an actual RSA-3072 identity, created
the bundled point-vector layer, performed masking and hexagonal binning, exposed
the Security \& Sharing key-exchange journey, and completed a real UI-driven
signed AES-256-GCM encrypt--decrypt round trip.  The 24 screen states reproduced
in Figures~\ref{fig:mapsafe-workflow-entry}--\ref{fig:mapsafe-blockchain-boundary}
were extracted from these instrumented workflows.  Thus, the screenshots are
not design mock-ups; they are outputs of the implemented Android application.
 when overriding the default, e.g.:
%   \newcommand{\mapsafefigdir}{paper/mapsafe-results/screenshots}
\providecommand{\mapsafefigdir}{screenshots}

\section{Results}
\label{sec:mapsafe-results}

The implementation produced an integrated MapSafe Mobile workflow within the
existing NextGIS-based Android application.  The resulting interface supports
two complementary modes.  The guided \emph{Protect \& Share} path presents
spatial representation before encryption, whereas the individual-feature path
allows anonymisation, encryption, and the future notarisation stage to be
entered separately.  Access operations are separated into verification,
decryption, and layer viewing.  Figure~\ref{fig:mapsafe-workflow-entry} shows
these entry points and their relationship.  All nested dialogs and activities
provide a labelled Back control and Android system Back follows the same parent
hierarchy; thus, for example, returning from \emph{Security \& Sharing} leads to
the MapSafe chooser rather than exiting the host map activity.

To make the workflow reproducible when no suitable field layer is available,
the chooser, guided workflow, and anonymisation screen provide a prominent
\emph{Use sample dataset} action (Fig.~\ref{fig:mapsafe-workflow-entry}a--b).
This action creates a known local point-vector layer containing 30 synthetic
observations in six clusters around Suva, selects it, and zooms the map to its
extent.  The sample is created directly as the point geometry type expected by
the anonymisation processors.  Consequently, it also provides a reliable path
past unsupported-file or ``not vector'' errors during demonstrations and
evaluation.

\begin{figure*}[t]
  \centering
  \begin{subfigure}[t]{0.235\textwidth}
    \includegraphics[width=\linewidth]{\mapsafefigdir/paper-01-mapsafe-workflow-chooser.png}
    \caption{MapSafe workflow chooser and sample-data path.}
    \label{fig:mapsafe-workflow-chooser}
  \end{subfigure}\hfill
  \begin{subfigure}[t]{0.235\textwidth}
    \includegraphics[width=\linewidth]{\mapsafefigdir/paper-02-guided-protect-share.png}
    \caption{Guided representation-before-encryption path.}
    \label{fig:mapsafe-guided-workflow}
  \end{subfigure}\hfill
  \begin{subfigure}[t]{0.235\textwidth}
    \includegraphics[width=\linewidth]{\mapsafefigdir/paper-03-safeguard-features.png}
    \caption{Individually selectable safeguard features.}
    \label{fig:mapsafe-safeguard-menu}
  \end{subfigure}\hfill
  \begin{subfigure}[t]{0.235\textwidth}
    \includegraphics[width=\linewidth]{\mapsafefigdir/paper-07-access-features.png}
    \caption{Verification, decryption, and viewing entry points.}
    \label{fig:mapsafe-access-menu}
  \end{subfigure}
  \caption{MapSafe Mobile workflow organisation.  The user may follow the
  ordered Protect \& Share sequence or invoke safeguards individually.  The
  full MapSafe wordmark identifies the main chooser, while compact headers use
  the icon-only mark.}
  \label{fig:mapsafe-workflow-entry}
\end{figure*}

\subsection{Anonymisation and spatial representation}
\label{subsec:mapsafe-anonymisation-results}

The anonymisation branch offers halo (donut) masking and hexagonal binning
(Fig.~\ref{fig:mapsafe-anonymisation-controls}).  Halo masking preserves every
feature and its non-spatial attributes but replaces each WGS84 point
$p_i=(\lambda_i,\phi_i)$ with a masked point $q_i$.  For each feature, a
cryptographically secure random generator selects a bearing uniformly from
$[0,2\pi)$ and a displacement uniformly from the user-selected interval
$[r_{\min},r_{\max})$.  The destination is calculated on a sphere rather than
by adding degree offsets.  The source layer remains unchanged and the candidate
masked layer is displayed in blue, allowing it to be compared with the precise
red reference points.

The implementation calculates Spruill's disclosure-risk measure immediately
after every masking attempt.  For $n$ paired original and masked points, define

\begin{equation}
 I_i =
 \begin{cases}
 1, & d(q_i,p_i)=\min_{1\leq j\leq n} d(q_i,p_j),\\
 0, & \text{otherwise},
 \end{cases}
 \qquad
 R_{\mathrm{Spruill}}=\frac{100}{n}\sum_{i=1}^{n}I_i .
 \label{eq:spruill-risk}
\end{equation}

Here, a point contributes to disclosure risk when the closest original location
to its masked position is still its own parent.  Nearest-neighbour searches are
performed with a three-dimensional unit-sphere $k$-d tree; unit-sphere chord
distance is monotonic with great-circle distance and avoids latitude-dependent
planar distortion.  Equal-distance ties that include the parent are counted as
parent-nearest.  Because a higher value is easier to interpret as stronger
privacy, the interface reports the inverted measure

\begin{equation}
 P_{\mathrm{inv}} = 100-R_{\mathrm{Spruill}}.
 \label{eq:spruill-inverted}
\end{equation}

In the illustrative 30-point run captured in
Figure~\ref{fig:mapsafe-anonymisation-result}a, 4 points remained
parent-nearest.  This gave $R_{\mathrm{Spruill}}=13.33\%$ and
$P_{\mathrm{inv}}=86.67/100$, with a mean displacement of 1041.21~m for the
selected 100--2000~m halo.  These values describe this stochastic run rather
than a fixed benchmark result.  The result screen exposes both the aggregate
score and its component counts.  The user may accept the candidate or select
\emph{Remask}; importantly, remasking returns to the same controls but starts
again from the precise source layer, never from the already displaced
candidate (Fig.~\ref{fig:mapsafe-anonymisation-result}b).  This prevents
unintended compounding of random offsets.

Hexagonal binning removes the individual point representation.  Each point is
assigned to a cell at the selected resolution, points with the same cell ID are
grouped, and one polygon containing a count attribute is written for each
occupied cell.  Native ARM and ARM64 builds use H3.  The x86/x86\_64 emulator
used for these captures lacks the H3 native library and therefore selected the
native-free Web Mercator hexagonal grid, whose resolution-dependent cell sizes
approximate the corresponding H3 choices.  The protected map in
Figure~\ref{fig:mapsafe-anonymisation-result}c demonstrates the resulting
masked points and counted polygon representation over the Suva sample.

\begin{figure*}[t]
  \centering
  \begin{subfigure}[t]{0.31\textwidth}
    \includegraphics[width=\linewidth]{\mapsafefigdir/paper-04-anonymise-options.png}
    \caption{Choice of masking or aggregation.}
    \label{fig:mapsafe-anonymise-choice}
  \end{subfigure}\hfill
  \begin{subfigure}[t]{0.31\textwidth}
    \includegraphics[width=\linewidth]{\mapsafefigdir/paper-16-halo-configuration.png}
    \caption{Halo bounds and privacy control.}
    \label{fig:mapsafe-halo-controls}
  \end{subfigure}\hfill
  \begin{subfigure}[t]{0.31\textwidth}
    \includegraphics[width=\linewidth]{\mapsafefigdir/paper-19-hexbin-configuration.png}
    \caption{Hexagon resolution and count aggregation.}
    \label{fig:mapsafe-hexbin-controls}
  \end{subfigure}
  \caption{Anonymisation controls.  Halo masking retains one displaced feature
  per input feature, whereas hexagonal binning replaces locations with occupied
  cells and their counts.}
  \label{fig:mapsafe-anonymisation-controls}
\end{figure*}

\begin{figure*}[t]
  \centering
  \begin{subfigure}[t]{0.31\textwidth}
    \includegraphics[width=\linewidth]{\mapsafefigdir/paper-17-halo-spruill-result.png}
    \caption{Inverted Spruill score and component measures.}
    \label{fig:mapsafe-spruill-result}
  \end{subfigure}\hfill
  \begin{subfigure}[t]{0.31\textwidth}
    \includegraphics[width=\linewidth]{\mapsafefigdir/paper-18-halo-remask.png}
    \caption{Remasking explicitly restarts from the precise source.}
    \label{fig:mapsafe-remask}
  \end{subfigure}\hfill
  \begin{subfigure}[t]{0.31\textwidth}
    \includegraphics[width=\linewidth]{\mapsafefigdir/paper-20-protected-map.png}
    \caption{Map output with precise references, masked points, and hexagons.}
    \label{fig:mapsafe-protected-map}
  \end{subfigure}
  \caption{Anonymisation outcomes for the bundled Suva point dataset.  The
  numeric privacy result is shown before the candidate is accepted, preserving
  an opportunity to remask or change parameters.}
  \label{fig:mapsafe-anonymisation-result}
\end{figure*}

\subsection{Encryption identity and public-key exchange}
\label{subsec:mapsafe-key-results}

The public-key workflow is surfaced through \emph{Security \& Sharing}
(Fig.~\ref{fig:mapsafe-key-exchange}).  Its detailed sequence is as follows.

\begin{enumerate}
  \item The user creates a local OpenPGP identity from a name, organisation,
  optional email address, and a passphrase of at least 12 characters
  (Fig.~\ref{fig:mapsafe-key-exchange}a).  The device generates an RSA-3072
  primary key for certification and signing plus a separate RSA-3072 encryption
  subkey.  The resulting full fingerprint is displayed to the user
  (Fig.~\ref{fig:mapsafe-key-exchange}b).

  \item The secret key material is encrypted with the user's passphrase and is
  additionally wrapped at rest by a non-exportable AES key in Android Keystore.
  It is not uploaded during ordinary key exchange.  Only the public key is
  exportable or publishable; an explicitly protected private-key backup is
  available as a recovery operation.

  \item The user selects an authenticated NextGIS account and a trusted NextGIS
  group (Fig.~\ref{fig:mapsafe-key-exchange}c).  MapSafe publishes the local
  public key to a group-specific key directory and can download keys belonging
  to current group members (Fig.~\ref{fig:mapsafe-key-exchange}d).  Directory
  ownership and current authentication-group membership are checked during
  synchronisation.

  \item Downloaded keys are initially discovered, not silently trusted.  The
  interface instructs the user to compare the complete fingerprint through an
  independent trusted channel and then accept it explicitly.  Acceptance pins
  that fingerprint to the member.  A later key change enters a pending/quarantined
  state, while changed, missing, duplicate, revoked, or removed-member keys are
  blocked from new recipient selection.  Previously accepted public keys remain
  cached so that encryption for established recipients can proceed offline.
\end{enumerate}

This design uses the authenticated group to establish membership and distribute
public material, but does not equate group membership with cryptographic key
authenticity.  The explicit fingerprint step supplies that missing binding and
prevents an updated directory entry from silently replacing a recipient's key.

\begin{figure*}[t]
  \centering
  \begin{subfigure}[t]{0.235\textwidth}
    \includegraphics[width=\linewidth]{\mapsafefigdir/paper-10-identity-creation.png}
    \caption{Local identity and passphrase entry.}
  \end{subfigure}\hfill
  \begin{subfigure}[t]{0.235\textwidth}
    \includegraphics[width=\linewidth]{\mapsafefigdir/paper-11-identity-success.png}
    \caption{Generated identity, key ID, and fingerprint.}
  \end{subfigure}\hfill
  \begin{subfigure}[t]{0.235\textwidth}
    \includegraphics[width=\linewidth]{\mapsafefigdir/paper-14-security-account-group.png}
    \caption{NextGIS account and group selection.}
  \end{subfigure}\hfill
  \begin{subfigure}[t]{0.235\textwidth}
    \includegraphics[width=\linewidth]{\mapsafefigdir/paper-15-security-identity-keys.png}
    \caption{Publishing, synchronising, and reviewing group keys.}
  \end{subfigure}
  \caption{Encryption identity and public-key exchange.  Private key material
  remains locally protected; public keys are distributed within a selected
  NextGIS group and become eligible recipients only after fingerprint review.}
  \label{fig:mapsafe-key-exchange}
\end{figure*}

\subsection{Multi-recipient encryption and signing}
\label{subsec:mapsafe-encryption-results}

Encryption may receive the currently selected protected map layer directly or
an existing file from device storage (Fig.~\ref{fig:mapsafe-encryption-flow}a).
For a map layer, MapSafe first exports the active representation as WGS84
GeoJSON, so the precise input need not leave the device.  The encryption screen
then shows the local identity, accepted recipient count, and the selected
content-protection profile (Fig.~\ref{fig:mapsafe-encryption-flow}b).

The package uses the RFC~9580 version-2 Symmetrically Encrypted Integrity
Protected Data profile with AES-256-GCM.  A fresh random content-encryption key
encrypts the compressed file payload once.  For each checked recipient, OpenPGP
adds a public-key encrypted session-key packet: the recipient's public
encryption subkey protects the content key, not a second copy of the entire
dataset.  Therefore adding recipients increases only the key-envelope portion
of the package.  The device owner's own public key can be selected as a
self-recipient, ensuring that the sender can subsequently decrypt the package.
The selection dialog makes this relationship explicit
(Fig.~\ref{fig:mapsafe-encryption-flow}c).

When \emph{Sign with my local identity} is selected, the passphrase unlocks the
local signing key only for the operation and a SHA-256 OpenPGP binary-document
signature is embedded with the payload.  Encryption provides confidentiality
and AEAD integrity, whereas this signature separately allows the receiver to
test which key created the package.  The successful state reports the
AES-256-GCM profile, authorised-recipient count, and whether a signed package
was saved (Fig.~\ref{fig:mapsafe-encryption-flow}d).

\begin{figure*}[t]
  \centering
  \begin{subfigure}[t]{0.235\textwidth}
    \includegraphics[width=\linewidth]{\mapsafefigdir/paper-05-encrypt-source-chooser.png}
    \caption{Protected-layer or file source selection.}
  \end{subfigure}\hfill
  \begin{subfigure}[t]{0.235\textwidth}
    \includegraphics[width=\linewidth]{\mapsafefigdir/paper-21-encrypt-protect.png}
    \caption{Encryption identity and protection summary.}
  \end{subfigure}\hfill
  \begin{subfigure}[t]{0.235\textwidth}
    \includegraphics[width=\linewidth]{\mapsafefigdir/paper-12-recipient-selection.png}
    \caption{Verified recipients and optional signing.}
  \end{subfigure}\hfill
  \begin{subfigure}[t]{0.235\textwidth}
    \includegraphics[width=\linewidth]{\mapsafefigdir/paper-22-encryption-success.png}
    \caption{Signed package ready to share.}
  \end{subfigure}
  \caption{Multi-recipient encryption flow.  The payload is encrypted once with
  AES-256-GCM; a separately public-key-encrypted session key grants access to
  each selected recipient.}
  \label{fig:mapsafe-encryption-flow}
\end{figure*}

\subsection{Decryption, verification, and viewing}
\label{subsec:mapsafe-decryption-results}

The access workflow starts with an OpenPGP package selected from storage
(Fig.~\ref{fig:mapsafe-decryption-flow}a--b).  MapSafe inspects the package's
recipient identifiers and chooses a matching locally held encryption secret
key.  The user passphrase unlocks that key; the resulting private key decrypts
the recipient's session-key packet, after which AES-256-GCM decrypts the single
payload.  A package addressed only to other recipients is rejected before
plaintext is released.

Two independent checks follow.  First, the SEIPD/AEAD integrity result must be
valid; an unprotected or modified package is rejected.  Second, if an embedded
signature is present, MapSafe locates the corresponding verification public key
and reports the signature as valid, invalid, or from an unknown signer.  The
decrypted file is held as a candidate until these checks complete.  In the
tested path the interface reported both a passed integrity check and a
cryptographically valid signature before offering \emph{Import and open}
(Fig.~\ref{fig:mapsafe-decryption-flow}c).  The final screen retains this status
as visible evidence (Fig.~\ref{fig:mapsafe-decryption-flow}d).

For a valid GeoJSON result, choosing \emph{Import and open} creates a local
NextGIS vector layer, selects it, and zooms to its extent.  Invalidly signed data
is not offered for map import.  This final promotion step separates successful
cryptographic processing from viewing and prevents an unauthenticated temporary
output from being mistaken for an accepted layer.

\begin{figure*}[t]
  \centering
  \begin{subfigure}[t]{0.235\textwidth}
    \includegraphics[width=\linewidth]{\mapsafefigdir/paper-08-decrypt-package-chooser.png}
    \caption{OpenPGP package entry point.}
  \end{subfigure}\hfill
  \begin{subfigure}[t]{0.235\textwidth}
    \includegraphics[width=\linewidth]{\mapsafefigdir/paper-23-decrypt-access.png}
    \caption{Local private-key and integrity summary.}
  \end{subfigure}\hfill
  \begin{subfigure}[t]{0.235\textwidth}
    \includegraphics[width=\linewidth]{\mapsafefigdir/paper-13-decryption-integrity-prompt.png}
    \caption{Verified candidate offered for import.}
  \end{subfigure}\hfill
  \begin{subfigure}[t]{0.235\textwidth}
    \includegraphics[width=\linewidth]{\mapsafefigdir/paper-24-decryption-success.png}
    \caption{Persistent success and signature status.}
  \end{subfigure}
  \caption{Decryption and access flow.  MapSafe first proves recipient access,
  then verifies AEAD integrity and the optional OpenPGP signature, and only then
  offers a valid spatial payload for import and viewing.}
  \label{fig:mapsafe-decryption-flow}
\end{figure*}

\subsection{Notarisation and blockchain-verification boundary}
\label{subsec:mapsafe-blockchain-results}

The proposed notarisation design records a SHA-256 digest of the encrypted
package plus minimal metadata, while the encrypted dataset remains off-chain.
Verification would independently hash the candidate package and compare that
digest with the value associated with a blockchain receipt; equality would
confirm byte-level correspondence with the notarised artifact, while inequality
would identify modification or selection of the wrong package.  This check is
complementary to OpenPGP: blockchain comparison provides timestamped external
evidence for a package digest, whereas decryption, AEAD integrity, and signature
verification establish authorised access and package authenticity.

However, the present Android build does \emph{not} contain a production
blockchain client and did not submit a transaction.  The interface deliberately
states that the production network is not configured and disables the minting
action (Fig.~\ref{fig:mapsafe-blockchain-boundary}a); the verification menu is
likewise marked as a later-stage interface
(Fig.~\ref{fig:mapsafe-blockchain-boundary}b).  This is an important result
boundary: these screens demonstrate workflow placement and privacy constraints,
not completed on-chain notarisation.  Automated scenarios currently exercise
the SHA-256 receipt comparison with an in-memory test ledger, including rejection
of a modified artifact, but those tests are not evidence of deployment to a
real network.

\begin{figure*}[t]
  \centering
  \begin{subfigure}[t]{0.31\textwidth}
    \includegraphics[width=\linewidth]{\mapsafefigdir/paper-06-notarisation-placeholder.png}
    \caption{Notarisation screen with the production action disabled.}
  \end{subfigure}\hspace{0.05\textwidth}
  \begin{subfigure}[t]{0.31\textwidth}
    \includegraphics[width=\linewidth]{\mapsafefigdir/paper-09-verification-placeholder.png}
    \caption{Blockchain-verification interface identified as future work.}
  \end{subfigure}
  \caption{Current blockchain boundary.  No real blockchain transaction is
  simulated or presented as completed by this build.}
  \label{fig:mapsafe-blockchain-boundary}
\end{figure*}

\subsection{Automated validation}
\label{subsec:mapsafe-validation-results}

The host-side MapSafe test suite executed 27 tests with no failures or skipped
tests.  These comprised eight OpenPGP engine tests, six public-key trust-state
tests, five Spruill-measure tests, two donut-masking tests, two portable-hex-grid
tests, and four end-to-end workflow-scenario tests.  Their coverage included
multi-recipient encryption, signed and unsigned packages, integrity/tamper
rejection, recipient mismatch, fingerprint changes, workflow combinations, and
the in-memory digest ledger.

Five targeted instrumented tests were then executed on a Pixel~9a Android
Virtual Device, again with no failures or skipped tests.  They traversed and
captured the workflow hierarchy, generated an actual RSA-3072 identity, created
the bundled point-vector layer, performed masking and hexagonal binning, exposed
the Security \& Sharing key-exchange journey, and completed a real UI-driven
signed AES-256-GCM encrypt--decrypt round trip.  The 24 screen states reproduced
in Figures~\ref{fig:mapsafe-workflow-entry}--\ref{fig:mapsafe-blockchain-boundary}
were extracted from these instrumented workflows.  Thus, the screenshots are
not design mock-ups; they are outputs of the implemented Android application.
 when overriding the default, e.g.:
%   \newcommand{\mapsafefigdir}{paper/mapsafe-results/screenshots}
\providecommand{\mapsafefigdir}{screenshots}

\section{Results}
\label{sec:mapsafe-results}

The implementation produced an integrated MapSafe Mobile workflow within the
existing NextGIS-based Android application.  The resulting interface supports
two complementary modes.  The guided \emph{Protect \& Share} path presents
spatial representation before encryption, whereas the individual-feature path
allows anonymisation, encryption, and the future notarisation stage to be
entered separately.  Access operations are separated into verification,
decryption, and layer viewing.  Figure~\ref{fig:mapsafe-workflow-entry} shows
these entry points and their relationship.  All nested dialogs and activities
provide a labelled Back control and Android system Back follows the same parent
hierarchy; thus, for example, returning from \emph{Security \& Sharing} leads to
the MapSafe chooser rather than exiting the host map activity.

To make the workflow reproducible when no suitable field layer is available,
the chooser, guided workflow, and anonymisation screen provide a prominent
\emph{Use sample dataset} action (Fig.~\ref{fig:mapsafe-workflow-entry}a--b).
This action creates a known local point-vector layer containing 30 synthetic
observations in six clusters around Suva, selects it, and zooms the map to its
extent.  The sample is created directly as the point geometry type expected by
the anonymisation processors.  Consequently, it also provides a reliable path
past unsupported-file or ``not vector'' errors during demonstrations and
evaluation.

\begin{figure*}[t]
  \centering
  \begin{subfigure}[t]{0.235\textwidth}
    \includegraphics[width=\linewidth]{\mapsafefigdir/paper-01-mapsafe-workflow-chooser.png}
    \caption{MapSafe workflow chooser and sample-data path.}
    \label{fig:mapsafe-workflow-chooser}
  \end{subfigure}\hfill
  \begin{subfigure}[t]{0.235\textwidth}
    \includegraphics[width=\linewidth]{\mapsafefigdir/paper-02-guided-protect-share.png}
    \caption{Guided representation-before-encryption path.}
    \label{fig:mapsafe-guided-workflow}
  \end{subfigure}\hfill
  \begin{subfigure}[t]{0.235\textwidth}
    \includegraphics[width=\linewidth]{\mapsafefigdir/paper-03-safeguard-features.png}
    \caption{Individually selectable safeguard features.}
    \label{fig:mapsafe-safeguard-menu}
  \end{subfigure}\hfill
  \begin{subfigure}[t]{0.235\textwidth}
    \includegraphics[width=\linewidth]{\mapsafefigdir/paper-07-access-features.png}
    \caption{Verification, decryption, and viewing entry points.}
    \label{fig:mapsafe-access-menu}
  \end{subfigure}
  \caption{MapSafe Mobile workflow organisation.  The user may follow the
  ordered Protect \& Share sequence or invoke safeguards individually.  The
  full MapSafe wordmark identifies the main chooser, while compact headers use
  the icon-only mark.}
  \label{fig:mapsafe-workflow-entry}
\end{figure*}

\subsection{Anonymisation and spatial representation}
\label{subsec:mapsafe-anonymisation-results}

The anonymisation branch offers halo (donut) masking and hexagonal binning
(Fig.~\ref{fig:mapsafe-anonymisation-controls}).  Halo masking preserves every
feature and its non-spatial attributes but replaces each WGS84 point
$p_i=(\lambda_i,\phi_i)$ with a masked point $q_i$.  For each feature, a
cryptographically secure random generator selects a bearing uniformly from
$[0,2\pi)$ and a displacement uniformly from the user-selected interval
$[r_{\min},r_{\max})$.  The destination is calculated on a sphere rather than
by adding degree offsets.  The source layer remains unchanged and the candidate
masked layer is displayed in blue, allowing it to be compared with the precise
red reference points.

The implementation calculates Spruill's disclosure-risk measure immediately
after every masking attempt.  For $n$ paired original and masked points, define

\begin{equation}
 I_i =
 \begin{cases}
 1, & d(q_i,p_i)=\min_{1\leq j\leq n} d(q_i,p_j),\\
 0, & \text{otherwise},
 \end{cases}
 \qquad
 R_{\mathrm{Spruill}}=\frac{100}{n}\sum_{i=1}^{n}I_i .
 \label{eq:spruill-risk}
\end{equation}

Here, a point contributes to disclosure risk when the closest original location
to its masked position is still its own parent.  Nearest-neighbour searches are
performed with a three-dimensional unit-sphere $k$-d tree; unit-sphere chord
distance is monotonic with great-circle distance and avoids latitude-dependent
planar distortion.  Equal-distance ties that include the parent are counted as
parent-nearest.  Because a higher value is easier to interpret as stronger
privacy, the interface reports the inverted measure

\begin{equation}
 P_{\mathrm{inv}} = 100-R_{\mathrm{Spruill}}.
 \label{eq:spruill-inverted}
\end{equation}

In the illustrative 30-point run captured in
Figure~\ref{fig:mapsafe-anonymisation-result}a, 4 points remained
parent-nearest.  This gave $R_{\mathrm{Spruill}}=13.33\%$ and
$P_{\mathrm{inv}}=86.67/100$, with a mean displacement of 1041.21~m for the
selected 100--2000~m halo.  These values describe this stochastic run rather
than a fixed benchmark result.  The result screen exposes both the aggregate
score and its component counts.  The user may accept the candidate or select
\emph{Remask}; importantly, remasking returns to the same controls but starts
again from the precise source layer, never from the already displaced
candidate (Fig.~\ref{fig:mapsafe-anonymisation-result}b).  This prevents
unintended compounding of random offsets.

Hexagonal binning removes the individual point representation.  Each point is
assigned to a cell at the selected resolution, points with the same cell ID are
grouped, and one polygon containing a count attribute is written for each
occupied cell.  Native ARM and ARM64 builds use H3.  The x86/x86\_64 emulator
used for these captures lacks the H3 native library and therefore selected the
native-free Web Mercator hexagonal grid, whose resolution-dependent cell sizes
approximate the corresponding H3 choices.  The protected map in
Figure~\ref{fig:mapsafe-anonymisation-result}c demonstrates the resulting
masked points and counted polygon representation over the Suva sample.

\begin{figure*}[t]
  \centering
  \begin{subfigure}[t]{0.31\textwidth}
    \includegraphics[width=\linewidth]{\mapsafefigdir/paper-04-anonymise-options.png}
    \caption{Choice of masking or aggregation.}
    \label{fig:mapsafe-anonymise-choice}
  \end{subfigure}\hfill
  \begin{subfigure}[t]{0.31\textwidth}
    \includegraphics[width=\linewidth]{\mapsafefigdir/paper-16-halo-configuration.png}
    \caption{Halo bounds and privacy control.}
    \label{fig:mapsafe-halo-controls}
  \end{subfigure}\hfill
  \begin{subfigure}[t]{0.31\textwidth}
    \includegraphics[width=\linewidth]{\mapsafefigdir/paper-19-hexbin-configuration.png}
    \caption{Hexagon resolution and count aggregation.}
    \label{fig:mapsafe-hexbin-controls}
  \end{subfigure}
  \caption{Anonymisation controls.  Halo masking retains one displaced feature
  per input feature, whereas hexagonal binning replaces locations with occupied
  cells and their counts.}
  \label{fig:mapsafe-anonymisation-controls}
\end{figure*}

\begin{figure*}[t]
  \centering
  \begin{subfigure}[t]{0.31\textwidth}
    \includegraphics[width=\linewidth]{\mapsafefigdir/paper-17-halo-spruill-result.png}
    \caption{Inverted Spruill score and component measures.}
    \label{fig:mapsafe-spruill-result}
  \end{subfigure}\hfill
  \begin{subfigure}[t]{0.31\textwidth}
    \includegraphics[width=\linewidth]{\mapsafefigdir/paper-18-halo-remask.png}
    \caption{Remasking explicitly restarts from the precise source.}
    \label{fig:mapsafe-remask}
  \end{subfigure}\hfill
  \begin{subfigure}[t]{0.31\textwidth}
    \includegraphics[width=\linewidth]{\mapsafefigdir/paper-20-protected-map.png}
    \caption{Map output with precise references, masked points, and hexagons.}
    \label{fig:mapsafe-protected-map}
  \end{subfigure}
  \caption{Anonymisation outcomes for the bundled Suva point dataset.  The
  numeric privacy result is shown before the candidate is accepted, preserving
  an opportunity to remask or change parameters.}
  \label{fig:mapsafe-anonymisation-result}
\end{figure*}

\subsection{Encryption identity and public-key exchange}
\label{subsec:mapsafe-key-results}

The public-key workflow is surfaced through \emph{Security \& Sharing}
(Fig.~\ref{fig:mapsafe-key-exchange}).  Its detailed sequence is as follows.

\begin{enumerate}
  \item The user creates a local OpenPGP identity from a name, organisation,
  optional email address, and a passphrase of at least 12 characters
  (Fig.~\ref{fig:mapsafe-key-exchange}a).  The device generates an RSA-3072
  primary key for certification and signing plus a separate RSA-3072 encryption
  subkey.  The resulting full fingerprint is displayed to the user
  (Fig.~\ref{fig:mapsafe-key-exchange}b).

  \item The secret key material is encrypted with the user's passphrase and is
  additionally wrapped at rest by a non-exportable AES key in Android Keystore.
  It is not uploaded during ordinary key exchange.  Only the public key is
  exportable or publishable; an explicitly protected private-key backup is
  available as a recovery operation.

  \item The user selects an authenticated NextGIS account and a trusted NextGIS
  group (Fig.~\ref{fig:mapsafe-key-exchange}c).  MapSafe publishes the local
  public key to a group-specific key directory and can download keys belonging
  to current group members (Fig.~\ref{fig:mapsafe-key-exchange}d).  Directory
  ownership and current authentication-group membership are checked during
  synchronisation.

  \item Downloaded keys are initially discovered, not silently trusted.  The
  interface instructs the user to compare the complete fingerprint through an
  independent trusted channel and then accept it explicitly.  Acceptance pins
  that fingerprint to the member.  A later key change enters a pending/quarantined
  state, while changed, missing, duplicate, revoked, or removed-member keys are
  blocked from new recipient selection.  Previously accepted public keys remain
  cached so that encryption for established recipients can proceed offline.
\end{enumerate}

This design uses the authenticated group to establish membership and distribute
public material, but does not equate group membership with cryptographic key
authenticity.  The explicit fingerprint step supplies that missing binding and
prevents an updated directory entry from silently replacing a recipient's key.

\begin{figure*}[t]
  \centering
  \begin{subfigure}[t]{0.235\textwidth}
    \includegraphics[width=\linewidth]{\mapsafefigdir/paper-10-identity-creation.png}
    \caption{Local identity and passphrase entry.}
  \end{subfigure}\hfill
  \begin{subfigure}[t]{0.235\textwidth}
    \includegraphics[width=\linewidth]{\mapsafefigdir/paper-11-identity-success.png}
    \caption{Generated identity, key ID, and fingerprint.}
  \end{subfigure}\hfill
  \begin{subfigure}[t]{0.235\textwidth}
    \includegraphics[width=\linewidth]{\mapsafefigdir/paper-14-security-account-group.png}
    \caption{NextGIS account and group selection.}
  \end{subfigure}\hfill
  \begin{subfigure}[t]{0.235\textwidth}
    \includegraphics[width=\linewidth]{\mapsafefigdir/paper-15-security-identity-keys.png}
    \caption{Publishing, synchronising, and reviewing group keys.}
  \end{subfigure}
  \caption{Encryption identity and public-key exchange.  Private key material
  remains locally protected; public keys are distributed within a selected
  NextGIS group and become eligible recipients only after fingerprint review.}
  \label{fig:mapsafe-key-exchange}
\end{figure*}

\subsection{Multi-recipient encryption and signing}
\label{subsec:mapsafe-encryption-results}

Encryption may receive the currently selected protected map layer directly or
an existing file from device storage (Fig.~\ref{fig:mapsafe-encryption-flow}a).
For a map layer, MapSafe first exports the active representation as WGS84
GeoJSON, so the precise input need not leave the device.  The encryption screen
then shows the local identity, accepted recipient count, and the selected
content-protection profile (Fig.~\ref{fig:mapsafe-encryption-flow}b).

The package uses the RFC~9580 version-2 Symmetrically Encrypted Integrity
Protected Data profile with AES-256-GCM.  A fresh random content-encryption key
encrypts the compressed file payload once.  For each checked recipient, OpenPGP
adds a public-key encrypted session-key packet: the recipient's public
encryption subkey protects the content key, not a second copy of the entire
dataset.  Therefore adding recipients increases only the key-envelope portion
of the package.  The device owner's own public key can be selected as a
self-recipient, ensuring that the sender can subsequently decrypt the package.
The selection dialog makes this relationship explicit
(Fig.~\ref{fig:mapsafe-encryption-flow}c).

When \emph{Sign with my local identity} is selected, the passphrase unlocks the
local signing key only for the operation and a SHA-256 OpenPGP binary-document
signature is embedded with the payload.  Encryption provides confidentiality
and AEAD integrity, whereas this signature separately allows the receiver to
test which key created the package.  The successful state reports the
AES-256-GCM profile, authorised-recipient count, and whether a signed package
was saved (Fig.~\ref{fig:mapsafe-encryption-flow}d).

\begin{figure*}[t]
  \centering
  \begin{subfigure}[t]{0.235\textwidth}
    \includegraphics[width=\linewidth]{\mapsafefigdir/paper-05-encrypt-source-chooser.png}
    \caption{Protected-layer or file source selection.}
  \end{subfigure}\hfill
  \begin{subfigure}[t]{0.235\textwidth}
    \includegraphics[width=\linewidth]{\mapsafefigdir/paper-21-encrypt-protect.png}
    \caption{Encryption identity and protection summary.}
  \end{subfigure}\hfill
  \begin{subfigure}[t]{0.235\textwidth}
    \includegraphics[width=\linewidth]{\mapsafefigdir/paper-12-recipient-selection.png}
    \caption{Verified recipients and optional signing.}
  \end{subfigure}\hfill
  \begin{subfigure}[t]{0.235\textwidth}
    \includegraphics[width=\linewidth]{\mapsafefigdir/paper-22-encryption-success.png}
    \caption{Signed package ready to share.}
  \end{subfigure}
  \caption{Multi-recipient encryption flow.  The payload is encrypted once with
  AES-256-GCM; a separately public-key-encrypted session key grants access to
  each selected recipient.}
  \label{fig:mapsafe-encryption-flow}
\end{figure*}

\subsection{Decryption, verification, and viewing}
\label{subsec:mapsafe-decryption-results}

The access workflow starts with an OpenPGP package selected from storage
(Fig.~\ref{fig:mapsafe-decryption-flow}a--b).  MapSafe inspects the package's
recipient identifiers and chooses a matching locally held encryption secret
key.  The user passphrase unlocks that key; the resulting private key decrypts
the recipient's session-key packet, after which AES-256-GCM decrypts the single
payload.  A package addressed only to other recipients is rejected before
plaintext is released.

Two independent checks follow.  First, the SEIPD/AEAD integrity result must be
valid; an unprotected or modified package is rejected.  Second, if an embedded
signature is present, MapSafe locates the corresponding verification public key
and reports the signature as valid, invalid, or from an unknown signer.  The
decrypted file is held as a candidate until these checks complete.  In the
tested path the interface reported both a passed integrity check and a
cryptographically valid signature before offering \emph{Import and open}
(Fig.~\ref{fig:mapsafe-decryption-flow}c).  The final screen retains this status
as visible evidence (Fig.~\ref{fig:mapsafe-decryption-flow}d).

For a valid GeoJSON result, choosing \emph{Import and open} creates a local
NextGIS vector layer, selects it, and zooms to its extent.  Invalidly signed data
is not offered for map import.  This final promotion step separates successful
cryptographic processing from viewing and prevents an unauthenticated temporary
output from being mistaken for an accepted layer.

\begin{figure*}[t]
  \centering
  \begin{subfigure}[t]{0.235\textwidth}
    \includegraphics[width=\linewidth]{\mapsafefigdir/paper-08-decrypt-package-chooser.png}
    \caption{OpenPGP package entry point.}
  \end{subfigure}\hfill
  \begin{subfigure}[t]{0.235\textwidth}
    \includegraphics[width=\linewidth]{\mapsafefigdir/paper-23-decrypt-access.png}
    \caption{Local private-key and integrity summary.}
  \end{subfigure}\hfill
  \begin{subfigure}[t]{0.235\textwidth}
    \includegraphics[width=\linewidth]{\mapsafefigdir/paper-13-decryption-integrity-prompt.png}
    \caption{Verified candidate offered for import.}
  \end{subfigure}\hfill
  \begin{subfigure}[t]{0.235\textwidth}
    \includegraphics[width=\linewidth]{\mapsafefigdir/paper-24-decryption-success.png}
    \caption{Persistent success and signature status.}
  \end{subfigure}
  \caption{Decryption and access flow.  MapSafe first proves recipient access,
  then verifies AEAD integrity and the optional OpenPGP signature, and only then
  offers a valid spatial payload for import and viewing.}
  \label{fig:mapsafe-decryption-flow}
\end{figure*}

\subsection{Notarisation and blockchain-verification boundary}
\label{subsec:mapsafe-blockchain-results}

The proposed notarisation design records a SHA-256 digest of the encrypted
package plus minimal metadata, while the encrypted dataset remains off-chain.
Verification would independently hash the candidate package and compare that
digest with the value associated with a blockchain receipt; equality would
confirm byte-level correspondence with the notarised artifact, while inequality
would identify modification or selection of the wrong package.  This check is
complementary to OpenPGP: blockchain comparison provides timestamped external
evidence for a package digest, whereas decryption, AEAD integrity, and signature
verification establish authorised access and package authenticity.

However, the present Android build does \emph{not} contain a production
blockchain client and did not submit a transaction.  The interface deliberately
states that the production network is not configured and disables the minting
action (Fig.~\ref{fig:mapsafe-blockchain-boundary}a); the verification menu is
likewise marked as a later-stage interface
(Fig.~\ref{fig:mapsafe-blockchain-boundary}b).  This is an important result
boundary: these screens demonstrate workflow placement and privacy constraints,
not completed on-chain notarisation.  Automated scenarios currently exercise
the SHA-256 receipt comparison with an in-memory test ledger, including rejection
of a modified artifact, but those tests are not evidence of deployment to a
real network.

\begin{figure*}[t]
  \centering
  \begin{subfigure}[t]{0.31\textwidth}
    \includegraphics[width=\linewidth]{\mapsafefigdir/paper-06-notarisation-placeholder.png}
    \caption{Notarisation screen with the production action disabled.}
  \end{subfigure}\hspace{0.05\textwidth}
  \begin{subfigure}[t]{0.31\textwidth}
    \includegraphics[width=\linewidth]{\mapsafefigdir/paper-09-verification-placeholder.png}
    \caption{Blockchain-verification interface identified as future work.}
  \end{subfigure}
  \caption{Current blockchain boundary.  No real blockchain transaction is
  simulated or presented as completed by this build.}
  \label{fig:mapsafe-blockchain-boundary}
\end{figure*}

\subsection{Automated validation}
\label{subsec:mapsafe-validation-results}

The host-side MapSafe test suite executed 27 tests with no failures or skipped
tests.  These comprised eight OpenPGP engine tests, six public-key trust-state
tests, five Spruill-measure tests, two donut-masking tests, two portable-hex-grid
tests, and four end-to-end workflow-scenario tests.  Their coverage included
multi-recipient encryption, signed and unsigned packages, integrity/tamper
rejection, recipient mismatch, fingerprint changes, workflow combinations, and
the in-memory digest ledger.

Five targeted instrumented tests were then executed on a Pixel~9a Android
Virtual Device, again with no failures or skipped tests.  They traversed and
captured the workflow hierarchy, generated an actual RSA-3072 identity, created
the bundled point-vector layer, performed masking and hexagonal binning, exposed
the Security \& Sharing key-exchange journey, and completed a real UI-driven
signed AES-256-GCM encrypt--decrypt round trip.  The 24 screen states reproduced
in Figures~\ref{fig:mapsafe-workflow-entry}--\ref{fig:mapsafe-blockchain-boundary}
were extracted from these instrumented workflows.  Thus, the screenshots are
not design mock-ups; they are outputs of the implemented Android application.
